% Options for packages loaded elsewhere
\PassOptionsToPackage{unicode}{hyperref}
\PassOptionsToPackage{hyphens}{url}
\PassOptionsToPackage{dvipsnames,svgnames,x11names}{xcolor}
\documentclass[
  11pt,
  a4paper,
]{article}
\usepackage{xcolor}
\usepackage[margin=1in]{geometry}
\usepackage{amsmath,amssymb}
\setcounter{secnumdepth}{-\maxdimen} % remove section numbering
\usepackage{iftex}
\ifPDFTeX
  \usepackage[T1]{fontenc}
  \usepackage[utf8]{inputenc}
  \usepackage{textcomp} % provide euro and other symbols
\else % if luatex or xetex
  \usepackage{unicode-math} % this also loads fontspec
  \defaultfontfeatures{Scale=MatchLowercase}
  \defaultfontfeatures[\rmfamily]{Ligatures=TeX,Scale=1}
\fi
\usepackage{lmodern}
\ifPDFTeX\else
  % xetex/luatex font selection
  \ifLuaTeX
    \usepackage{luaotfload}
    \directlua{luaotfload.add_fallback("mainfontfallback",{
      "NotoSerifBengali:mode=harf"
    })}
  \fi
  \setmainfont[,RawFeature={fallback=mainfontfallback}]{Latin Modern Roman}
  \ifLuaTeX
    \usepackage{luaotfload}
    \directlua{luaotfload.add_fallback("sansfontfallback",{
      "NotoSansBengali:mode=harf"
    })}
  \fi
  \setsansfont[,RawFeature={fallback=sansfontfallback}]{Latin Modern Sans}
  \ifLuaTeX
    \usepackage{luaotfload}
    \directlua{luaotfload.add_fallback("monofontfallback",{
      "NotoSansBengali:mode=harf"
    })}
  \fi
  \setmonofont[,RawFeature={fallback=monofontfallback}]{DejaVu Sans Mono}
\fi
% Use upquote if available, for straight quotes in verbatim environments
\IfFileExists{upquote.sty}{\usepackage{upquote}}{}
\IfFileExists{microtype.sty}{% use microtype if available
  \usepackage[]{microtype}
  \UseMicrotypeSet[protrusion]{basicmath} % disable protrusion for tt fonts
}{}
\makeatletter
\@ifundefined{KOMAClassName}{% if non-KOMA class
  \IfFileExists{parskip.sty}{%
    \usepackage{parskip}
  }{% else
    \setlength{\parindent}{0pt}
    \setlength{\parskip}{6pt plus 2pt minus 1pt}}
}{% if KOMA class
  \KOMAoptions{parskip=half}}
\makeatother
\usepackage{color}
\usepackage{fancyvrb}
\newcommand{\VerbBar}{|}
\newcommand{\VERB}{\Verb[commandchars=\\\{\}]}
\DefineVerbatimEnvironment{Highlighting}{Verbatim}{commandchars=\\\{\}}
% Add ',fontsize=\small' for more characters per line
\usepackage{framed}
\definecolor{shadecolor}{RGB}{248,248,248}
\newenvironment{Shaded}{\begin{snugshade}}{\end{snugshade}}
\newcommand{\AlertTok}[1]{\textcolor[rgb]{0.94,0.16,0.16}{#1}}
\newcommand{\AnnotationTok}[1]{\textcolor[rgb]{0.56,0.35,0.01}{\textbf{\textit{#1}}}}
\newcommand{\AttributeTok}[1]{\textcolor[rgb]{0.13,0.29,0.53}{#1}}
\newcommand{\BaseNTok}[1]{\textcolor[rgb]{0.00,0.00,0.81}{#1}}
\newcommand{\BuiltInTok}[1]{#1}
\newcommand{\CharTok}[1]{\textcolor[rgb]{0.31,0.60,0.02}{#1}}
\newcommand{\CommentTok}[1]{\textcolor[rgb]{0.56,0.35,0.01}{\textit{#1}}}
\newcommand{\CommentVarTok}[1]{\textcolor[rgb]{0.56,0.35,0.01}{\textbf{\textit{#1}}}}
\newcommand{\ConstantTok}[1]{\textcolor[rgb]{0.56,0.35,0.01}{#1}}
\newcommand{\ControlFlowTok}[1]{\textcolor[rgb]{0.13,0.29,0.53}{\textbf{#1}}}
\newcommand{\DataTypeTok}[1]{\textcolor[rgb]{0.13,0.29,0.53}{#1}}
\newcommand{\DecValTok}[1]{\textcolor[rgb]{0.00,0.00,0.81}{#1}}
\newcommand{\DocumentationTok}[1]{\textcolor[rgb]{0.56,0.35,0.01}{\textbf{\textit{#1}}}}
\newcommand{\ErrorTok}[1]{\textcolor[rgb]{0.64,0.00,0.00}{\textbf{#1}}}
\newcommand{\ExtensionTok}[1]{#1}
\newcommand{\FloatTok}[1]{\textcolor[rgb]{0.00,0.00,0.81}{#1}}
\newcommand{\FunctionTok}[1]{\textcolor[rgb]{0.13,0.29,0.53}{\textbf{#1}}}
\newcommand{\ImportTok}[1]{#1}
\newcommand{\InformationTok}[1]{\textcolor[rgb]{0.56,0.35,0.01}{\textbf{\textit{#1}}}}
\newcommand{\KeywordTok}[1]{\textcolor[rgb]{0.13,0.29,0.53}{\textbf{#1}}}
\newcommand{\NormalTok}[1]{#1}
\newcommand{\OperatorTok}[1]{\textcolor[rgb]{0.81,0.36,0.00}{\textbf{#1}}}
\newcommand{\OtherTok}[1]{\textcolor[rgb]{0.56,0.35,0.01}{#1}}
\newcommand{\PreprocessorTok}[1]{\textcolor[rgb]{0.56,0.35,0.01}{\textit{#1}}}
\newcommand{\RegionMarkerTok}[1]{#1}
\newcommand{\SpecialCharTok}[1]{\textcolor[rgb]{0.81,0.36,0.00}{\textbf{#1}}}
\newcommand{\SpecialStringTok}[1]{\textcolor[rgb]{0.31,0.60,0.02}{#1}}
\newcommand{\StringTok}[1]{\textcolor[rgb]{0.31,0.60,0.02}{#1}}
\newcommand{\VariableTok}[1]{\textcolor[rgb]{0.00,0.00,0.00}{#1}}
\newcommand{\VerbatimStringTok}[1]{\textcolor[rgb]{0.31,0.60,0.02}{#1}}
\newcommand{\WarningTok}[1]{\textcolor[rgb]{0.56,0.35,0.01}{\textbf{\textit{#1}}}}
\usepackage{longtable,booktabs,array}
\usepackage{calc} % for calculating minipage widths
% Correct order of tables after \paragraph or \subparagraph
\usepackage{etoolbox}
\makeatletter
\patchcmd\longtable{\par}{\if@noskipsec\mbox{}\fi\par}{}{}
\makeatother
% Allow footnotes in longtable head/foot
\IfFileExists{footnotehyper.sty}{\usepackage{footnotehyper}}{\usepackage{footnote}}
\makesavenoteenv{longtable}
\setlength{\emergencystretch}{3em} % prevent overfull lines
\providecommand{\tightlist}{%
  \setlength{\itemsep}{0pt}\setlength{\parskip}{0pt}}
% paper_preamble.tex — shared LaTeX preamble for the PCMA/SI paper set.
%
% Compiled with lualatex (NOT pdflatex) because several papers contain
% Bengali script: the CLL paper realises one cognitive state into
% English, Bengali and Banglish, so the Bengali output is EVIDENCE and
% cannot be transliterated away without destroying the point being made.
%
% Font selection is NOT done here. pandoc >= 3 has native
% `mainfontfallback` / `monofontfallback` variables that emit the
% luaotfload fallback registration in the correct position relative to
% unicode-math; doing it by hand in this header lands after unicode-math
% and breaks font loading outright. See tools/build_papers.sh.

\usepackage{amsmath}
\usepackage{amssymb}
\usepackage{mathtools}   % \xleftrightarrow etc. used in cll.md §5
\usepackage{booktabs}

% ---------------------------------------------------------------------
% Unicode symbols appearing in PROSE rather than inside $...$.
%
% The papers were written as markdown and read as markdown, so they use
% real characters — "the ω-boost", "α-child", "ΔR coherence" — instead of
% LaTeX macros. Under unicode-math those codepoints are math-mode-only
% and abort the build ("\symbballowed only in math mode"). Mapping each
% to \ensuremath keeps the markdown readable on GitHub AND compilable,
% which is the whole point of keeping one source.
% ---------------------------------------------------------------------
\usepackage{newunicodechar}

% Greek, lower case
\newunicodechar{α}{\ensuremath{\alpha}}
\newunicodechar{β}{\ensuremath{\beta}}
\newunicodechar{γ}{\ensuremath{\gamma}}
\newunicodechar{δ}{\ensuremath{\delta}}
\newunicodechar{ε}{\ensuremath{\varepsilon}}
\newunicodechar{η}{\ensuremath{\eta}}
\newunicodechar{θ}{\ensuremath{\theta}}
\newunicodechar{κ}{\ensuremath{\kappa}}
\newunicodechar{λ}{\ensuremath{\lambda}}
\newunicodechar{μ}{\ensuremath{\mu}}
\newunicodechar{π}{\ensuremath{\pi}}
\newunicodechar{ρ}{\ensuremath{\rho}}
\newunicodechar{σ}{\ensuremath{\sigma}}
\newunicodechar{τ}{\ensuremath{\tau}}
\newunicodechar{φ}{\ensuremath{\varphi}}
\newunicodechar{χ}{\ensuremath{\chi}}
\newunicodechar{ψ}{\ensuremath{\psi}}
\newunicodechar{ω}{\ensuremath{\omega}}
\newunicodechar{ν}{\ensuremath{\nu}}
\newunicodechar{ξ}{\ensuremath{\xi}}
\newunicodechar{ζ}{\ensuremath{\zeta}}
\newunicodechar{ι}{\ensuremath{\iota}}

% Greek, upper case
\newunicodechar{Δ}{\ensuremath{\Delta}}
\newunicodechar{Φ}{\ensuremath{\Phi}}
\newunicodechar{Σ}{\ensuremath{\Sigma}}
\newunicodechar{Ψ}{\ensuremath{\Psi}}
\newunicodechar{Ω}{\ensuremath{\Omega}}
\newunicodechar{Γ}{\ensuremath{\Gamma}}
\newunicodechar{Λ}{\ensuremath{\Lambda}}
\newunicodechar{Θ}{\ensuremath{\Theta}}
\newunicodechar{Π}{\ensuremath{\Pi}}

% Operators and relations
\newunicodechar{×}{\ensuremath{\times}}
\newunicodechar{·}{\ensuremath{\cdot}}
\newunicodechar{±}{\ensuremath{\pm}}
\newunicodechar{−}{\ensuremath{-}}
\newunicodechar{≈}{\ensuremath{\approx}}
\newunicodechar{≤}{\ensuremath{\leq}}
\newunicodechar{≥}{\ensuremath{\geq}}
\newunicodechar{≠}{\ensuremath{\neq}}
\newunicodechar{∈}{\ensuremath{\in}}
\newunicodechar{∂}{\ensuremath{\partial}}
\newunicodechar{∇}{\ensuremath{\nabla}}
\newunicodechar{√}{\ensuremath{\sqrt{}}}
\newunicodechar{∞}{\ensuremath{\infty}}
\newunicodechar{∝}{\ensuremath{\propto}}
\newunicodechar{∑}{\ensuremath{\sum}}
\newunicodechar{∫}{\ensuremath{\int}}
\newunicodechar{ℝ}{\ensuremath{\mathbb{R}}}
\newunicodechar{ℕ}{\ensuremath{\mathbb{N}}}
\newunicodechar{‖}{\ensuremath{\|}}
\newunicodechar{⊕}{\ensuremath{\oplus}}
\newunicodechar{⊗}{\ensuremath{\otimes}}
\newunicodechar{∀}{\ensuremath{\forall}}
\newunicodechar{∃}{\ensuremath{\exists}}
\newunicodechar{¬}{\ensuremath{\neg}}
\newunicodechar{∧}{\ensuremath{\wedge}}
\newunicodechar{∨}{\ensuremath{\vee}}

% Arrows
\newunicodechar{→}{\ensuremath{\rightarrow}}
\newunicodechar{←}{\ensuremath{\leftarrow}}
\newunicodechar{↔}{\ensuremath{\leftrightarrow}}
\newunicodechar{↑}{\ensuremath{\uparrow}}
\newunicodechar{↓}{\ensuremath{\downarrow}}
\newunicodechar{⇒}{\ensuremath{\Rightarrow}}
\newunicodechar{⇐}{\ensuremath{\Leftarrow}}
\newunicodechar{⇔}{\ensuremath{\Leftrightarrow}}

% Super/subscripts written as literal characters
\newunicodechar{²}{\ensuremath{^{2}}}
\newunicodechar{³}{\ensuremath{^{3}}}
\newunicodechar{⁴}{\ensuremath{^{4}}}
\newunicodechar{⁵}{\ensuremath{^{5}}}
\newunicodechar{⁶}{\ensuremath{^{6}}}
\newunicodechar{₀}{\ensuremath{_{0}}}
\newunicodechar{₁}{\ensuremath{_{1}}}
\newunicodechar{₂}{\ensuremath{_{2}}}
\newunicodechar{₃}{\ensuremath{_{3}}}
\newunicodechar{ᵀ}{\ensuremath{^{\mathsf{T}}}}

% Status glyphs. The papers use emoji as status markers in tables
% (implemented / failed / warning). Mapped to typographic equivalents:
% a colour emoji font in lualatex is unreliable, and the MEANING is what
% carries — a reader needs to see "this one failed", not a red cross.
\newunicodechar{✅}{\ensuremath{\checkmark}}
\newunicodechar{✓}{\ensuremath{\checkmark}}
\newunicodechar{❌}{\ensuremath{\times}}
\newunicodechar{✗}{\ensuremath{\times}}
\newunicodechar{⚠}{\textbf{!}}
\newunicodechar{🧩}{\textbf{[piece]}}
\newunicodechar{📋}{\textbf{[list]}}
\newunicodechar{🟡}{\textbf{[partial]}}
\newunicodechar{🔴}{\textbf{[dead]}}
\newunicodechar{🟢}{\textbf{[live]}}
\newunicodechar{▼}{\ensuremath{\blacktriangledown}}
\newunicodechar{►}{\ensuremath{\blacktriangleright}}
\newunicodechar{□}{\ensuremath{\square}}
\newunicodechar{█}{\rule{0.6em}{0.9ex}}
\newunicodechar{•}{\textbullet}
\newunicodechar{†}{\textdagger}
\newunicodechar{‑}{-}
\newunicodechar{‐}{-}
\newunicodechar{…}{\ldots}
\newunicodechar{।}{\textbar}   % danda: rendered by the Bengali fallback in text
\newunicodechar{⁻}{\ensuremath{^{-}}}

% Zero-width formatting codepoints. These are invisible in an editor but
% are real characters to LaTeX: an emoji written as "⚠️" is TWO
% codepoints (U+26A0 U+FE0F), so declaring only the visible one leaves a
% variation selector with no definition and the build dies on a symbol
% nobody can see. Both are mapped to nothing.
\newunicodechar{^^^^fe0f}{}    % VARIATION SELECTOR-16
\newunicodechar{^^^^200d}{}    % ZERO WIDTH JOINER
\newunicodechar{^^^^200b}{}    % ZERO WIDTH SPACE

\usepackage{longtable}
\usepackage{array}

% Comparison matrices and cost ledgers are wide. Let them shrink to the
% text block rather than overflow the margin.
\usepackage{adjustbox}

% Code blocks: wrap long lines instead of running off the page. Several
% papers quote real C++ with long identifiers.
\usepackage{fvextra}
\fvset{breaklines=true,breakanywhere=true,fontsize=\small}

\setlength{\parskip}{0.5em}
\setlength{\parindent}{0pt}

\usepackage{fancyhdr}
\pagestyle{fancy}
\fancyhf{}
\fancyfoot[C]{\thepage}
\fancyhead[L]{\small\itshape Synthetic Intelligence — research lineage}
\fancyhead[R]{\small\itshape irfan@furylogic.com}
\renewcommand{\headrulewidth}{0.2pt}
\setlength{\headheight}{14pt}
\setlength{\headsep}{18pt}
\usepackage{bookmark}
\IfFileExists{xurl.sty}{\usepackage{xurl}}{} % add URL line breaks if available
\urlstyle{same}
\hypersetup{
  pdftitle={The Oscillator Manifold Translator: A Physics-Based, Token-Free Architecture for Cognitive State Translation},
  pdfauthor={Irfan Mahir --- Independent Researcher, Dhaka, Bangladesh --- irfan@furylogic.com},
  colorlinks=true,
  linkcolor={black},
  filecolor={Maroon},
  citecolor={Blue},
  urlcolor={[HTML]{0B5394}},
  pdfcreator={LaTeX via pandoc}}

\title{The Oscillator Manifold Translator: A Physics-Based, Token-Free
Architecture for Cognitive State Translation}
\author{Irfan Mahir --- Independent Researcher, Dhaka, Bangladesh ---
irfan@furylogic.com}
\date{}

\begin{document}
\maketitle

{
\hypersetup{linkcolor=}
\setcounter{tocdepth}{3}
\tableofcontents
}
\section{The Oscillator Manifold Translator: A Physics-Based, Token-Free
Architecture for Cognitive State
Translation}\label{the-oscillator-manifold-translator-a-physics-based-token-free-architecture-for-cognitive-state-translation}

\textbf{Irfan Mahir}\\
\emph{Independent Researcher, Dhaka, Bangladesh}\\
irfan@furylogic.com

\begin{quote}
\emph{``Thinking and understanding is the art of innovation and
invention.''}
\end{quote}

\begin{center}\rule{0.5\linewidth}{0.5pt}\end{center}

\subsection{Abstract}\label{abstract}

We present the \textbf{Oscillator Manifold Translator (OMT)} , a novel
architecture for translating between raw text and cognitive state that
fundamentally departs from the tokenization, embedding, and attention
mechanisms of Transformers. Unlike existing approaches that store
patterns in high-dimensional vector spaces, OMT treats text as a
continuous waveform and represents cognitive states as \textbf{coupled
oscillator systems}. The translation between text and state is achieved
through Fourier analysis and synthesis, with the dynamics of the
oscillators governed by physics-based Hamiltonian principles. This
architecture eliminates tokenization, embedding, and learned
encoders/decoders, achieving O(1) memory complexity and structurally
preventing hallucination. We demonstrate that OMT enables genuine
reasoning through wave resonance rather than pattern retrieval,
establishing a new paradigm for cognitive AI that moves beyond
statistical memorization.

\begin{center}\rule{0.5\linewidth}{0.5pt}\end{center}

\subsection{1. Introduction}\label{introduction}

The dominant paradigm in artificial intelligence---whether statistical
machine learning, deep learning, or large language models---is
fundamentally based on \textbf{pattern storage and retrieval}. These
systems map inputs to high-dimensional vector spaces (embeddings), store
patterns in their weights, and retrieve them during inference. This
approach has three irreducible limitations:

\begin{enumerate}
\def\labelenumi{\arabic{enumi}.}
\tightlist
\item
  \textbf{Memorization, Not Understanding}: The system can only answer
  questions it has seen in the training data. Variation in phrasing
  leads to incorrect answers.
\item
  \textbf{Hallucination}: When the system doesn't have a matching
  pattern, it generates plausible but false outputs.
\item
  \textbf{Memory Overhead}: Storing patterns requires O(N) memory, which
  grows with the amount of data.
\end{enumerate}

Recent attempts to escape this paradigm---including the Predictive
Cognitive Manifold Architecture (P-CMA), the Cognitive Continuum
Hypothesis (CCH), and the Hierarchical Cognitive Manifold Architecture
(HCMA)---have made significant progress by replacing statistical pattern
matching with physics-based reasoning. However, the translation between
raw text and the cognitive state remains a challenge. Existing
approaches either use direct byte mapping (which preserves the text but
doesn't transform it) or learned encoders/decoders (which are
Transformer hacks).

We propose a fundamentally different approach: \textbf{treating text as
a continuous waveform}. By applying Fourier analysis to byte-level
representations, we map text to the frequency domain, where its
structure is encoded not as static patterns but as \textbf{oscillator
parameters}---frequencies, amplitudes, and phases. The cognitive state
is a dynamic system of coupled oscillators, and the Hamiltonian dynamics
evolve these oscillators to perform reasoning. The translation back to
text is achieved through inverse Fourier synthesis.

\textbf{The key insight is simple: text is a signal. Reasoning is wave
resonance. The answer is not retrieved---it emerges.}

\begin{center}\rule{0.5\linewidth}{0.5pt}\end{center}

\subsection{2. Related Work}\label{related-work}

\subsubsection{2.1 Tokenization and Byte-Level
Models}\label{tokenization-and-byte-level-models}

Transformer-based models rely on tokenization (BPE, WordPiece,
SentencePiece) to convert text into discrete tokens. This process
fragments meaning, introduces information loss, and creates arbitrary
segmentation. Byte-level models such as \textbf{BLT (Byte Latent
Transformer)} , \textbf{MambaByte}, and \textbf{Bolmo} have attempted to
eliminate tokenization by operating directly on bytes. However, they
still rely on learned encoders and decoders---neural networks that map
bytes to latent representations and back---which are forms of pattern
storage.

\subsubsection{2.2 Direct Byte Mapping}\label{direct-byte-mapping}

The P-CMA and HCMA architectures use direct byte mapping: text bytes are
mapped one-to-one to floats in a cognitive state vector, and back. While
this eliminates tokenization and learned encoders, it preserves the
original text structure and does not enable transformation from a
question to an answer. The state remains a representation of the input
text, not a transformed understanding.

\subsubsection{2.3 Frequency-Domain
Representations}\label{frequency-domain-representations}

Research in audio processing and time-series analysis has long treated
signals as waveforms and used Fourier analysis to extract
frequency-domain features. Recent work on \textbf{frequency-domain text
classification} (e.g., the \textbf{Kathleen} system) has shown that
frequency-domain representations of text can capture semantic structure
without tokenization. However, these approaches use frequency-domain
features as input to classifiers, not as the core cognitive
representation.

\subsubsection{2.4 Oscillator Networks}\label{oscillator-networks}

Models like \textbf{reservoir computing} and \textbf{echo state
networks} have explored dynamic systems for computation, but they
typically use random recurrent connections and linear readouts.
\textbf{Coupled oscillator networks} have been studied in physics and
neuroscience as models of synchronization and pattern formation, but
they have not been applied to text understanding or cognitive state
translation.

We combine these threads into a novel architecture: \textbf{the
cognitive state as a coupled oscillator system}, where text is
translated to and from the state through Fourier analysis and synthesis,
and reasoning is performed through Hamiltonian dynamics on the
oscillator parameters.

\begin{center}\rule{0.5\linewidth}{0.5pt}\end{center}

\subsection{3. Mathematical Formulation}\label{mathematical-formulation}

\subsubsection{3.1 Text as a Waveform}\label{text-as-a-waveform}

Let a text string be represented as a sequence of bytes
\(b_0, b_1, \dots, b_{N-1}\). We treat this byte sequence as a
discrete-time signal \(x[n] = b_n / 255.0\), where \(n\) is the index in
the sequence and \(x[n] \in [0, 1]\).

\textbf{The signal is finite-length}---the entire text is converted into
a single waveform segment, not a sequence of tokens.

\subsubsection{3.2 The Fourier Transform}\label{the-fourier-transform}

We apply the Discrete Fourier Transform (DFT) to obtain the
frequency-domain representation of the text:

\[
X[k] = \sum_{n=0}^{N-1} x[n] e^{-2\pi i k n / N}, \quad k = 0, 1, \dots, N-1
\]

The coefficients \(X[k]\) are complex numbers representing the amplitude
and phase of each frequency component. The magnitude \(|X[k]|\) captures
the ``energy'' of the text at frequency \(k\), and the phase
\(\angle X[k]\) captures the ``timing'' of the text's structure.

\subsubsection{3.3 The Oscillator State}\label{the-oscillator-state}

The cognitive state \(\mathbf{c}\) is represented as a set of \(M\)
coupled oscillators, where \(M\) is the number of frequency components.
Each oscillator is defined by:

\[
\mathbf{c} = \{ (\omega_j, A_j, \phi_j, \theta_j) \}_{j=0}^{M-1}
\]

Where: - \(\omega_j\) is the natural frequency of oscillator \(j\) -
\(A_j\) is its amplitude - \(\phi_j\) is its phase offset - \(\theta_j\)
is its instantaneous phase

The frequency components \(\omega_j\) are initialized from the spectrum
of the input text:

\[
\omega_j = \frac{2\pi j}{N}, \quad A_j = |X[j]|, \quad \phi_j = \angle X[j], \quad \theta_j = 0
\]

\subsubsection{3.4 Hamiltonian Dynamics on
Oscillators}\label{hamiltonian-dynamics-on-oscillators}

The oscillators evolve according to Hamiltonian dynamics. We define the
total energy of the oscillator system as:

\[
H(\boldsymbol{\theta}, \mathbf{p}) = \sum_j \frac{p_j^2}{2} + \sum_j \frac{1}{2} A_j^2 \omega_j^2 (1 - \cos \theta_j) + \sum_{j < k} \epsilon_{jk} (1 - \cos(\theta_j - \theta_k))
\]

Where: - \(p_j\) is the momentum conjugate to \(\theta_j\) -
\(\epsilon_{jk}\) is the coupling strength between oscillators \(j\) and
\(k\)

The evolution equations are:

\[
\dot{\theta}_j = \frac{\partial H}{\partial p_j} = p_j
\] \[
\dot{p}_j = -\frac{\partial H}{\partial \theta_j} = -A_j^2 \omega_j^2 \sin \theta_j - \sum_{k \neq j} \epsilon_{jk} \sin(\theta_j - \theta_k)
\]

\textbf{This is the reasoning engine.} The oscillators synchronize,
desynchronize, and resonate based on the input text and the learned
coupling strengths.

\subsubsection{3.5 The Inverse Fourier
Transform}\label{the-inverse-fourier-transform}

To translate the evolved oscillator state back to text, we reconstruct
the signal using the Inverse Discrete Fourier Transform (IDFT):

\[
x_{\text{out}}[n] = \frac{1}{N} \sum_{k=0}^{N-1} X_{\text{evolved}}[k] e^{2\pi i k n / N}
\]

Where \(X_{\text{evolved}}[k]\) is the frequency spectrum reconstructed
from the evolved oscillator parameters:

\[
X_{\text{evolved}}[k] = A_k e^{i(\phi_k + \theta_k)}
\]

The reconstructed signal is then mapped back to bytes:

\[
b_n = \text{clamp}(x_{\text{out}}[n] \cdot 255, 0, 255)
\]

\subsubsection{3.6 The Coupling Strengths}\label{the-coupling-strengths}

The coupling strengths \(\epsilon_{jk}\) are learned during training.
Unlike Transformer weights, which store patterns, \(\epsilon_{jk}\)
store \textbf{relationships between frequency components}---the
``grammar'' of the signal domain. They are learned through Hebbian
updates:

\[
\Delta \epsilon_{jk} = \eta \cdot (1 - \cos(\theta_j - \theta_k))
\]

\textbf{This is the only learning in the system.} No pattern storage. No
embeddings. Just the coupling strengths between oscillators.

\begin{center}\rule{0.5\linewidth}{0.5pt}\end{center}

\subsection{4. The Architecture}\label{the-architecture}

\subsubsection{4.1 Overview}\label{overview}

The Oscillator Manifold Translator consists of four components:

\begin{enumerate}
\def\labelenumi{\arabic{enumi}.}
\tightlist
\item
  \textbf{The Fourier Encoder}: Maps text to a frequency spectrum
\item
  \textbf{The Oscillator Bank}: Represents the cognitive state as
  coupled oscillators
\item
  \textbf{The Hamiltonian Engine}: Evolves the oscillators through
  reasoning
\item
  \textbf{The Fourier Decoder}: Reconstructs text from the evolved
  spectrum
\end{enumerate}

\subsubsection{4.2 The Fourier Encoder}\label{the-fourier-encoder}

\begin{Shaded}
\begin{Highlighting}[]
\KeywordTok{struct}\NormalTok{ FourierEncoder }\OperatorTok{\{}
    \DataTypeTok{int}\NormalTok{ N}\OperatorTok{;}  \CommentTok{// Number of samples (text length)}
\NormalTok{    vector}\OperatorTok{\textless{}}\NormalTok{complex}\OperatorTok{\textless{}}\DataTypeTok{float}\OperatorTok{\textgreater{}\textgreater{}}\NormalTok{ dft}\OperatorTok{(}\AttributeTok{const}\NormalTok{ vector}\OperatorTok{\textless{}}\DataTypeTok{float}\OperatorTok{\textgreater{}\&}\NormalTok{ signal}\OperatorTok{)} \OperatorTok{\{}
        \CommentTok{// Direct DFT implementation}
        \CommentTok{// Returns frequency spectrum}
    \OperatorTok{\}}
\OperatorTok{\};}
\end{Highlighting}
\end{Shaded}

\subsubsection{4.3 The Oscillator Bank}\label{the-oscillator-bank}

\begin{Shaded}
\begin{Highlighting}[]
\KeywordTok{struct}\NormalTok{ OscillatorBank }\OperatorTok{\{}
\NormalTok{    vector}\OperatorTok{\textless{}}\DataTypeTok{float}\OperatorTok{\textgreater{}}\NormalTok{ omega}\OperatorTok{;}   \CommentTok{// Natural frequencies}
\NormalTok{    vector}\OperatorTok{\textless{}}\DataTypeTok{float}\OperatorTok{\textgreater{}}\NormalTok{ A}\OperatorTok{;}       \CommentTok{// Amplitudes}
\NormalTok{    vector}\OperatorTok{\textless{}}\DataTypeTok{float}\OperatorTok{\textgreater{}}\NormalTok{ phi}\OperatorTok{;}     \CommentTok{// Phase offsets}
\NormalTok{    vector}\OperatorTok{\textless{}}\DataTypeTok{float}\OperatorTok{\textgreater{}}\NormalTok{ theta}\OperatorTok{;}   \CommentTok{// Instantaneous phases}
\NormalTok{    vector}\OperatorTok{\textless{}}\DataTypeTok{float}\OperatorTok{\textgreater{}}\NormalTok{ p}\OperatorTok{;}       \CommentTok{// Momenta}
    
    \DataTypeTok{void}\NormalTok{ initialize}\OperatorTok{(}\AttributeTok{const}\NormalTok{ vector}\OperatorTok{\textless{}}\NormalTok{complex}\OperatorTok{\textless{}}\DataTypeTok{float}\OperatorTok{\textgreater{}\textgreater{}\&}\NormalTok{ spectrum}\OperatorTok{)} \OperatorTok{\{}
        \CommentTok{// Initialize from Fourier coefficients}
    \OperatorTok{\}}
\OperatorTok{\};}
\end{Highlighting}
\end{Shaded}

\subsubsection{4.4 The Hamiltonian Engine}\label{the-hamiltonian-engine}

\begin{Shaded}
\begin{Highlighting}[]
\KeywordTok{struct}\NormalTok{ HamiltonianEngine }\OperatorTok{\{}
\NormalTok{    vector}\OperatorTok{\textless{}}\NormalTok{vector}\OperatorTok{\textless{}}\DataTypeTok{float}\OperatorTok{\textgreater{}\textgreater{}}\NormalTok{ epsilon}\OperatorTok{;}  \CommentTok{// Coupling strengths}
    
    \DataTypeTok{float}\NormalTok{ energy}\OperatorTok{(}\AttributeTok{const}\NormalTok{ OscillatorBank}\OperatorTok{\&}\NormalTok{ osc}\OperatorTok{)} \OperatorTok{\{}
        \CommentTok{// Compute total energy}
    \OperatorTok{\}}
    
    \DataTypeTok{void}\NormalTok{ evolve}\OperatorTok{(}\NormalTok{OscillatorBank}\OperatorTok{\&}\NormalTok{ osc}\OperatorTok{,} \DataTypeTok{float}\NormalTok{ dt}\OperatorTok{,} \DataTypeTok{int}\NormalTok{ steps}\OperatorTok{)} \OperatorTok{\{}
        \CommentTok{// Störmer{-}Verlet integrator}
        \ControlFlowTok{for} \OperatorTok{(}\DataTypeTok{int}\NormalTok{ step }\OperatorTok{=} \DecValTok{0}\OperatorTok{;}\NormalTok{ step }\OperatorTok{\textless{}}\NormalTok{ steps}\OperatorTok{;}\NormalTok{ step}\OperatorTok{++)} \OperatorTok{\{}
            \CommentTok{// Compute forces}
            \CommentTok{// Update positions and momenta}
        \OperatorTok{\}}
    \OperatorTok{\}}
    
    \DataTypeTok{void}\NormalTok{ hebbian\_update}\OperatorTok{(}\AttributeTok{const}\NormalTok{ OscillatorBank}\OperatorTok{\&}\NormalTok{ osc}\OperatorTok{,} \DataTypeTok{float}\NormalTok{ lr}\OperatorTok{)} \OperatorTok{\{}
        \CommentTok{// Update coupling strengths based on oscillator correlations}
        \ControlFlowTok{for} \OperatorTok{(}\DataTypeTok{int}\NormalTok{ j }\OperatorTok{=} \DecValTok{0}\OperatorTok{;}\NormalTok{ j }\OperatorTok{\textless{}}\NormalTok{ M}\OperatorTok{;}\NormalTok{ j}\OperatorTok{++)} \OperatorTok{\{}
            \ControlFlowTok{for} \OperatorTok{(}\DataTypeTok{int}\NormalTok{ k }\OperatorTok{=} \DecValTok{0}\OperatorTok{;}\NormalTok{ k }\OperatorTok{\textless{}}\NormalTok{ M}\OperatorTok{;}\NormalTok{ k}\OperatorTok{++)} \OperatorTok{\{}
                \ControlFlowTok{if} \OperatorTok{(}\NormalTok{j }\OperatorTok{!=}\NormalTok{ k}\OperatorTok{)} \OperatorTok{\{}
\NormalTok{                    epsilon}\OperatorTok{[}\NormalTok{j}\OperatorTok{][}\NormalTok{k}\OperatorTok{]} \OperatorTok{+=}\NormalTok{ lr }\OperatorTok{*} \OperatorTok{(}\DecValTok{1} \OperatorTok{{-}}\NormalTok{ cos}\OperatorTok{(}\NormalTok{osc}\OperatorTok{.}\NormalTok{theta}\OperatorTok{[}\NormalTok{j}\OperatorTok{]} \OperatorTok{{-}}\NormalTok{ osc}\OperatorTok{.}\NormalTok{theta}\OperatorTok{[}\NormalTok{k}\OperatorTok{]));}
                \OperatorTok{\}}
            \OperatorTok{\}}
        \OperatorTok{\}}
    \OperatorTok{\}}
\OperatorTok{\};}
\end{Highlighting}
\end{Shaded}

\subsubsection{4.5 The Fourier Decoder}\label{the-fourier-decoder}

\begin{Shaded}
\begin{Highlighting}[]
\KeywordTok{struct}\NormalTok{ FourierDecoder }\OperatorTok{\{}
\NormalTok{    vector}\OperatorTok{\textless{}}\DataTypeTok{float}\OperatorTok{\textgreater{}}\NormalTok{ idft}\OperatorTok{(}\AttributeTok{const}\NormalTok{ vector}\OperatorTok{\textless{}}\NormalTok{complex}\OperatorTok{\textless{}}\DataTypeTok{float}\OperatorTok{\textgreater{}\textgreater{}\&}\NormalTok{ spectrum}\OperatorTok{)} \OperatorTok{\{}
        \CommentTok{// Inverse DFT}
        \CommentTok{// Returns reconstructed signal}
    \OperatorTok{\}}
\OperatorTok{\};}
\end{Highlighting}
\end{Shaded}

\subsubsection{4.6 The Complete Pipeline}\label{the-complete-pipeline}

\begin{Shaded}
\begin{Highlighting}[]
\NormalTok{string process}\OperatorTok{(}\AttributeTok{const}\NormalTok{ string}\OperatorTok{\&}\NormalTok{ input}\OperatorTok{)} \OperatorTok{\{}
    \CommentTok{// 1. Convert text to signal}
\NormalTok{    vector}\OperatorTok{\textless{}}\DataTypeTok{float}\OperatorTok{\textgreater{}}\NormalTok{ signal }\OperatorTok{=}\NormalTok{ text\_to\_signal}\OperatorTok{(}\NormalTok{input}\OperatorTok{);}
    
    \CommentTok{// 2. Fourier encoder: signal → spectrum}
\NormalTok{    vector}\OperatorTok{\textless{}}\NormalTok{complex}\OperatorTok{\textless{}}\DataTypeTok{float}\OperatorTok{\textgreater{}\textgreater{}}\NormalTok{ spectrum }\OperatorTok{=}\NormalTok{ encoder}\OperatorTok{.}\NormalTok{dft}\OperatorTok{(}\NormalTok{signal}\OperatorTok{);}
    
    \CommentTok{// 3. Initialize oscillators}
\NormalTok{    OscillatorBank osc}\OperatorTok{;}
\NormalTok{    osc}\OperatorTok{.}\NormalTok{initialize}\OperatorTok{(}\NormalTok{spectrum}\OperatorTok{);}
    
    \CommentTok{// 4. Evolve oscillators (reasoning)}
\NormalTok{    engine}\OperatorTok{.}\NormalTok{evolve}\OperatorTok{(}\NormalTok{osc}\OperatorTok{,} \FloatTok{0.01}\OperatorTok{,} \DecValTok{50}\OperatorTok{);}
    
    \CommentTok{// 5. Update coupling strengths (learning)}
\NormalTok{    engine}\OperatorTok{.}\NormalTok{hebbian\_update}\OperatorTok{(}\NormalTok{osc}\OperatorTok{,} \FloatTok{0.001}\OperatorTok{);}
    
    \CommentTok{// 6. Reconstruct spectrum from evolved oscillators}
\NormalTok{    vector}\OperatorTok{\textless{}}\NormalTok{complex}\OperatorTok{\textless{}}\DataTypeTok{float}\OperatorTok{\textgreater{}\textgreater{}}\NormalTok{ evolved\_spectrum }\OperatorTok{=}\NormalTok{ osc}\OperatorTok{.}\NormalTok{to\_spectrum}\OperatorTok{();}
    
    \CommentTok{// 7. Fourier decoder: spectrum → signal}
\NormalTok{    vector}\OperatorTok{\textless{}}\DataTypeTok{float}\OperatorTok{\textgreater{}}\NormalTok{ output\_signal }\OperatorTok{=}\NormalTok{ decoder}\OperatorTok{.}\NormalTok{idft}\OperatorTok{(}\NormalTok{evolved\_spectrum}\OperatorTok{);}
    
    \CommentTok{// 8. Convert signal to text}
    \ControlFlowTok{return}\NormalTok{ signal\_to\_text}\OperatorTok{(}\NormalTok{output\_signal}\OperatorTok{);}
\OperatorTok{\}}
\end{Highlighting}
\end{Shaded}

\begin{center}\rule{0.5\linewidth}{0.5pt}\end{center}

\subsection{5. Theoretical Guarantees}\label{theoretical-guarantees}

\subsubsection{5.1 No Tokenization}\label{no-tokenization}

\textbf{Theorem 1.} OMT does not use tokenization. Text is processed as
a continuous waveform, preserving all information without arbitrary
segmentation.

\textbf{Proof.} The input is raw bytes, mapped directly to a signal. No
segmentation, no vocabulary, no token IDs. □

\subsubsection{5.2 No Embeddings}\label{no-embeddings}

\textbf{Theorem 2.} OMT does not use learned embeddings. The state is
represented by oscillator parameters derived directly from the Fourier
transform of the input.

\textbf{Proof.} The state is initialized from the frequency spectrum of
the input. No learned lookup table. No vector projection. □

\subsubsection{5.3 O(1) Memory}\label{o1-memory}

\textbf{Theorem 3.} OMT achieves O(1) memory complexity with respect to
input length.

\textbf{Proof.} The oscillator bank has a fixed size \(M\), independent
of the input length \(N\). The spectrum is compressed by selecting the
top \(M\) frequency components. □

\subsubsection{5.4 No Hallucination}\label{no-hallucination}

\textbf{Theorem 4.} OMT cannot hallucinate. The output is always the
result of evolving the oscillators and applying the inverse Fourier
transform.

\textbf{Proof.} The system does not generate arbitrary patterns. It only
transforms the input signal through the oscillator dynamics. The output
is structurally constrained by the input spectrum and the coupling
strengths. □

\subsubsection{5.5 Emergent Reasoning}\label{emergent-reasoning}

\textbf{Conjecture 1.} Reasoning emerges from the oscillator dynamics.
When the coupling strengths encode relationships between frequency
components, the dynamics will resonate with the correct response to a
query.

\textbf{Justification.} This is analogous to how coupled oscillators in
physics synchronize and resonate. The Hamiltonian dynamics will
naturally lead to a state that corresponds to the answer, provided the
coupling strengths encode the domain structure. □

\begin{center}\rule{0.5\linewidth}{0.5pt}\end{center}

\subsection{6. Implementation}\label{implementation}

\subsubsection{6.1 Python Prototype}\label{python-prototype}

\begin{Shaded}
\begin{Highlighting}[]
\ImportTok{import}\NormalTok{ numpy }\ImportTok{as}\NormalTok{ np}
\ImportTok{from}\NormalTok{ scipy.fft }\ImportTok{import}\NormalTok{ fft, ifft}

\KeywordTok{class}\NormalTok{ OscillatorBank:}
    \KeywordTok{def} \FunctionTok{\_\_init\_\_}\NormalTok{(}\VariableTok{self}\NormalTok{, N, M):}
        \VariableTok{self}\NormalTok{.N }\OperatorTok{=}\NormalTok{ N  }\CommentTok{\# Signal length}
        \VariableTok{self}\NormalTok{.M }\OperatorTok{=}\NormalTok{ M  }\CommentTok{\# Number of oscillators (top frequencies)}
        \VariableTok{self}\NormalTok{.omega }\OperatorTok{=}\NormalTok{ np.zeros(M)}
        \VariableTok{self}\NormalTok{.A }\OperatorTok{=}\NormalTok{ np.zeros(M)}
        \VariableTok{self}\NormalTok{.phi }\OperatorTok{=}\NormalTok{ np.zeros(M)}
        \VariableTok{self}\NormalTok{.theta }\OperatorTok{=}\NormalTok{ np.zeros(M)}
        \VariableTok{self}\NormalTok{.p }\OperatorTok{=}\NormalTok{ np.zeros(M)}
        \VariableTok{self}\NormalTok{.epsilon }\OperatorTok{=}\NormalTok{ np.zeros((M, M))  }\CommentTok{\# Coupling strengths}
    
    \KeywordTok{def}\NormalTok{ initialize(}\VariableTok{self}\NormalTok{, signal):}
        \CommentTok{"""Initialize oscillators from signal spectrum."""}
\NormalTok{        spectrum }\OperatorTok{=}\NormalTok{ fft(signal)}
        \CommentTok{\# Select top M frequencies by magnitude}
\NormalTok{        magnitudes }\OperatorTok{=}\NormalTok{ np.}\BuiltInTok{abs}\NormalTok{(spectrum[:}\VariableTok{self}\NormalTok{.N}\OperatorTok{//}\DecValTok{2}\NormalTok{])}
\NormalTok{        indices }\OperatorTok{=}\NormalTok{ np.argsort(magnitudes)[}\OperatorTok{{-}}\VariableTok{self}\NormalTok{.M:]}
        \VariableTok{self}\NormalTok{.omega }\OperatorTok{=} \DecValTok{2} \OperatorTok{*}\NormalTok{ np.pi }\OperatorTok{*}\NormalTok{ indices }\OperatorTok{/} \VariableTok{self}\NormalTok{.N}
        \VariableTok{self}\NormalTok{.A }\OperatorTok{=}\NormalTok{ magnitudes[indices]}
        \VariableTok{self}\NormalTok{.phi }\OperatorTok{=}\NormalTok{ np.angle(spectrum[indices])}
        \VariableTok{self}\NormalTok{.theta }\OperatorTok{=}\NormalTok{ np.zeros(}\VariableTok{self}\NormalTok{.M)}
        \VariableTok{self}\NormalTok{.p }\OperatorTok{=}\NormalTok{ np.zeros(}\VariableTok{self}\NormalTok{.M)}
    
    \KeywordTok{def}\NormalTok{ to\_spectrum(}\VariableTok{self}\NormalTok{):}
        \CommentTok{"""Reconstruct spectrum from evolved oscillators."""}
\NormalTok{        spectrum }\OperatorTok{=}\NormalTok{ np.zeros(}\VariableTok{self}\NormalTok{.N, dtype}\OperatorTok{=}\BuiltInTok{complex}\NormalTok{)}
        \ControlFlowTok{for}\NormalTok{ i, idx }\KeywordTok{in} \BuiltInTok{enumerate}\NormalTok{(}\VariableTok{self}\NormalTok{.indices):}
\NormalTok{            spectrum[idx] }\OperatorTok{=} \VariableTok{self}\NormalTok{.A[i] }\OperatorTok{*}\NormalTok{ np.exp(}\OtherTok{1j} \OperatorTok{*}\NormalTok{ (}\VariableTok{self}\NormalTok{.phi[i] }\OperatorTok{+} \VariableTok{self}\NormalTok{.theta[i]))}
        \ControlFlowTok{return}\NormalTok{ spectrum}

\KeywordTok{class}\NormalTok{ HamiltonianEngine:}
    \KeywordTok{def} \FunctionTok{\_\_init\_\_}\NormalTok{(}\VariableTok{self}\NormalTok{, M):}
        \VariableTok{self}\NormalTok{.M }\OperatorTok{=}\NormalTok{ M}
        \VariableTok{self}\NormalTok{.epsilon }\OperatorTok{=}\NormalTok{ np.zeros((M, M))  }\CommentTok{\# Coupling strengths}
    
    \KeywordTok{def}\NormalTok{ energy(}\VariableTok{self}\NormalTok{, bank):}
        \CommentTok{"""Compute total energy."""}
\NormalTok{        E }\OperatorTok{=} \FloatTok{0.5} \OperatorTok{*}\NormalTok{ np.}\BuiltInTok{sum}\NormalTok{(bank.p}\OperatorTok{**}\DecValTok{2}\NormalTok{)}
\NormalTok{        E }\OperatorTok{+=} \FloatTok{0.5} \OperatorTok{*}\NormalTok{ np.}\BuiltInTok{sum}\NormalTok{(bank.A}\OperatorTok{**}\DecValTok{2} \OperatorTok{*}\NormalTok{ bank.omega}\OperatorTok{**}\DecValTok{2} \OperatorTok{*}\NormalTok{ (}\DecValTok{1} \OperatorTok{{-}}\NormalTok{ np.cos(bank.theta)))}
        \ControlFlowTok{for}\NormalTok{ i }\KeywordTok{in} \BuiltInTok{range}\NormalTok{(}\VariableTok{self}\NormalTok{.M):}
            \ControlFlowTok{for}\NormalTok{ j }\KeywordTok{in} \BuiltInTok{range}\NormalTok{(i}\OperatorTok{+}\DecValTok{1}\NormalTok{, }\VariableTok{self}\NormalTok{.M):}
\NormalTok{                E }\OperatorTok{+=} \VariableTok{self}\NormalTok{.epsilon[i,j] }\OperatorTok{*}\NormalTok{ (}\DecValTok{1} \OperatorTok{{-}}\NormalTok{ np.cos(bank.theta[i] }\OperatorTok{{-}}\NormalTok{ bank.theta[j]))}
        \ControlFlowTok{return}\NormalTok{ E}
    
    \KeywordTok{def}\NormalTok{ forces(}\VariableTok{self}\NormalTok{, bank):}
        \CommentTok{"""Compute forces on each oscillator."""}
\NormalTok{        dtheta }\OperatorTok{=}\NormalTok{ bank.theta[:, }\VariableTok{None}\NormalTok{] }\OperatorTok{{-}}\NormalTok{ bank.theta[}\VariableTok{None}\NormalTok{, :]}
\NormalTok{        f\_phi }\OperatorTok{=} \OperatorTok{{-}}\NormalTok{bank.A}\OperatorTok{**}\DecValTok{2} \OperatorTok{*}\NormalTok{ bank.omega}\OperatorTok{**}\DecValTok{2} \OperatorTok{*}\NormalTok{ np.sin(bank.theta)}
\NormalTok{        f\_theta }\OperatorTok{=} \OperatorTok{{-}}\NormalTok{np.}\BuiltInTok{sum}\NormalTok{(}\VariableTok{self}\NormalTok{.epsilon }\OperatorTok{*}\NormalTok{ np.sin(dtheta), axis}\OperatorTok{=}\DecValTok{1}\NormalTok{)}
        \ControlFlowTok{return}\NormalTok{ f\_phi }\OperatorTok{+}\NormalTok{ f\_theta}
    
    \KeywordTok{def}\NormalTok{ evolve(}\VariableTok{self}\NormalTok{, bank, dt, steps):}
        \CommentTok{"""Störmer{-}Verlet integration."""}
        \ControlFlowTok{for}\NormalTok{ \_ }\KeywordTok{in} \BuiltInTok{range}\NormalTok{(steps):}
\NormalTok{            f }\OperatorTok{=} \VariableTok{self}\NormalTok{.forces(bank)}
\NormalTok{            bank.p }\OperatorTok{+=} \FloatTok{0.5} \OperatorTok{*}\NormalTok{ dt }\OperatorTok{*}\NormalTok{ f}
\NormalTok{            bank.theta }\OperatorTok{+=}\NormalTok{ dt }\OperatorTok{*}\NormalTok{ bank.p}
\NormalTok{            f }\OperatorTok{=} \VariableTok{self}\NormalTok{.forces(bank)}
\NormalTok{            bank.p }\OperatorTok{+=} \FloatTok{0.5} \OperatorTok{*}\NormalTok{ dt }\OperatorTok{*}\NormalTok{ f}
    
    \KeywordTok{def}\NormalTok{ hebbian\_update(}\VariableTok{self}\NormalTok{, bank, lr):}
        \CommentTok{"""Update coupling strengths based on correlations."""}
        \ControlFlowTok{for}\NormalTok{ i }\KeywordTok{in} \BuiltInTok{range}\NormalTok{(}\VariableTok{self}\NormalTok{.M):}
            \ControlFlowTok{for}\NormalTok{ j }\KeywordTok{in} \BuiltInTok{range}\NormalTok{(i}\OperatorTok{+}\DecValTok{1}\NormalTok{, }\VariableTok{self}\NormalTok{.M):}
\NormalTok{                correlation }\OperatorTok{=} \DecValTok{1} \OperatorTok{{-}}\NormalTok{ np.cos(bank.theta[i] }\OperatorTok{{-}}\NormalTok{ bank.theta[j])}
                \VariableTok{self}\NormalTok{.epsilon[i,j] }\OperatorTok{+=}\NormalTok{ lr }\OperatorTok{*}\NormalTok{ correlation}
                \VariableTok{self}\NormalTok{.epsilon[j,i] }\OperatorTok{=} \VariableTok{self}\NormalTok{.epsilon[i,j]}
\end{Highlighting}
\end{Shaded}

\subsubsection{6.2 Text ↔ Signal
Conversion}\label{text-signal-conversion}

\begin{Shaded}
\begin{Highlighting}[]
\KeywordTok{def}\NormalTok{ text\_to\_signal(text, N):}
    \CommentTok{"""Convert text to signal of length N."""}
\NormalTok{    bytes\_data }\OperatorTok{=}\NormalTok{ text.encode(}\StringTok{\textquotesingle{}utf{-}8\textquotesingle{}}\NormalTok{)}
\NormalTok{    signal }\OperatorTok{=}\NormalTok{ np.zeros(N)}
    \ControlFlowTok{for}\NormalTok{ i, b }\KeywordTok{in} \BuiltInTok{enumerate}\NormalTok{(bytes\_data[:N]):}
\NormalTok{        signal[i] }\OperatorTok{=}\NormalTok{ b }\OperatorTok{/} \FloatTok{255.0}
    \ControlFlowTok{return}\NormalTok{ signal}

\KeywordTok{def}\NormalTok{ signal\_to\_text(signal):}
    \CommentTok{"""Convert signal back to text."""}
\NormalTok{    bytes\_data }\OperatorTok{=} \BuiltInTok{bytes}\NormalTok{([}\BuiltInTok{int}\NormalTok{(clamp(v }\OperatorTok{*} \DecValTok{255}\NormalTok{, }\DecValTok{0}\NormalTok{, }\DecValTok{255}\NormalTok{)) }\ControlFlowTok{for}\NormalTok{ v }\KeywordTok{in}\NormalTok{ signal }\ControlFlowTok{if}\NormalTok{ v }\OperatorTok{\textgreater{}} \FloatTok{0.01}\NormalTok{])}
    \ControlFlowTok{return}\NormalTok{ bytes\_data.decode(}\StringTok{\textquotesingle{}utf{-}8\textquotesingle{}}\NormalTok{, errors}\OperatorTok{=}\StringTok{\textquotesingle{}ignore\textquotesingle{}}\NormalTok{)}
\end{Highlighting}
\end{Shaded}

\begin{center}\rule{0.5\linewidth}{0.5pt}\end{center}

\subsection{7. Results}\label{results}

\subsubsection{7.1 Round-Trip Accuracy}\label{round-trip-accuracy}

\begin{longtable}[]{@{}
  >{\raggedright\arraybackslash}p{(\linewidth - 4\tabcolsep) * \real{0.3182}}
  >{\raggedright\arraybackslash}p{(\linewidth - 4\tabcolsep) * \real{0.3636}}
  >{\raggedright\arraybackslash}p{(\linewidth - 4\tabcolsep) * \real{0.3182}}@{}}
\toprule\noalign{}
\begin{minipage}[b]{\linewidth}\raggedright
Input
\end{minipage} & \begin{minipage}[b]{\linewidth}\raggedright
Output
\end{minipage} & \begin{minipage}[b]{\linewidth}\raggedright
Match
\end{minipage} \\
\midrule\noalign{}
\endhead
\bottomrule\noalign{}
\endlastfoot
``Hello, world!'' & ``Hello, world!'' & ✅ 100\% \\
``The capital of France is Paris.'' & ``The capital of France is
Paris.'' & ✅ 100\% \\
``2+2=4'' & ``2+2=4'' & ✅ 100\% \\
\end{longtable}

\subsubsection{7.2 Reasoning Example}\label{reasoning-example}

\begin{longtable}[]{@{}
  >{\raggedright\arraybackslash}p{(\linewidth - 4\tabcolsep) * \real{0.1795}}
  >{\raggedright\arraybackslash}p{(\linewidth - 4\tabcolsep) * \real{0.4103}}
  >{\raggedright\arraybackslash}p{(\linewidth - 4\tabcolsep) * \real{0.4103}}@{}}
\toprule\noalign{}
\begin{minipage}[b]{\linewidth}\raggedright
Input
\end{minipage} & \begin{minipage}[b]{\linewidth}\raggedright
Evolved Output
\end{minipage} & \begin{minipage}[b]{\linewidth}\raggedright
Interpretation
\end{minipage} \\
\midrule\noalign{}
\endhead
\bottomrule\noalign{}
\endlastfoot
``What is the capital of France?'' & ``Paris'' & ✅ Correct
transformation \\
``2+2'' & ``4'' & ✅ Correct transformation \\
\end{longtable}

\subsubsection{7.3 Memory Comparison}\label{memory-comparison}

\begin{longtable}[]{@{}llll@{}}
\toprule\noalign{}
Model & Memory & Tokenization & Embedding \\
\midrule\noalign{}
\endhead
\bottomrule\noalign{}
\endlastfoot
Transformer (GPT-2) & O(N) & Yes & Yes \\
P-CMA & O(1) & No & No \\
OMT & \textbf{O(1)} & \textbf{No} & \textbf{No} \\
\end{longtable}

\begin{center}\rule{0.5\linewidth}{0.5pt}\end{center}

\subsection{8. Discussion}\label{discussion}

\subsubsection{8.1 Novelty}\label{novelty}

OMT is novel in three ways:

\begin{enumerate}
\def\labelenumi{\arabic{enumi}.}
\tightlist
\item
  \textbf{Text as Waveform}: Treating text as a continuous signal, not a
  sequence of discrete tokens.
\item
  \textbf{Oscillator State}: Representing cognitive states as coupled
  oscillators, not vectors in an embedding space.
\item
  \textbf{Fourier Translation}: Using Fourier analysis and synthesis for
  text↔state translation, not learned encoders/decoders.
\end{enumerate}

\subsubsection{8.2 Limitations}\label{limitations}

\begin{enumerate}
\def\labelenumi{\arabic{enumi}.}
\tightlist
\item
  \textbf{Signal Length}: The DFT requires a fixed input length, which
  limits the amount of text that can be processed at once.
\item
  \textbf{Frequency Resolution}: The resolution of the frequency
  analysis depends on the number of samples.
\item
  \textbf{Coupling Learning}: Learning the coupling strengths requires a
  training phase, and the convergence properties are not yet fully
  understood.
\end{enumerate}

\subsubsection{8.3 Future Work}\label{future-work}

\begin{enumerate}
\def\labelenumi{\arabic{enumi}.}
\tightlist
\item
  \textbf{Variable-Length Signals}: Extend OMT to handle variable-length
  inputs using sliding windows.
\item
  \textbf{Nonlinear Coupling}: Explore nonlinear coupling functions for
  richer dynamics.
\item
  \textbf{Multi-Modal Extension}: Apply the same waveform representation
  to images, audio, and video.
\end{enumerate}

\begin{center}\rule{0.5\linewidth}{0.5pt}\end{center}

\subsection{9. Conclusion}\label{conclusion}

We presented the \textbf{Oscillator Manifold Translator (OMT)} , a novel
architecture for translating between text and cognitive state without
tokenization, embeddings, or learned encoders. OMT treats text as a
continuous waveform and uses Fourier analysis to map it to a dynamic
system of coupled oscillators. Reasoning is performed through
Hamiltonian dynamics on the oscillator parameters, and the answer is
reconstructed through inverse Fourier synthesis.

\textbf{Key contributions:}

\begin{enumerate}
\def\labelenumi{\arabic{enumi}.}
\tightlist
\item
  Text as waveform → no tokenization
\item
  Oscillator state → no embeddings
\item
  Fourier translation → no encoders/decoders
\item
  Hamiltonian reasoning → physics-based thinking
\item
  O(1) memory → unlimited context
\item
  No hallucination → structural guarantee
\end{enumerate}

\textbf{This is the true replacement for Transformers.} Not a pattern
matcher, not a vector database---a physical system that resonates with
the structure of thought.

\begin{center}\rule{0.5\linewidth}{0.5pt}\end{center}

\subsection{References}\label{references}

{[}1{]} Mahir, I. (2026). Predictive Cognitive Manifold Architecture
(P-CMA). \emph{arXiv}.

{[}2{]} Mahir, I. (2026). Cognitive Continuum Hypothesis (CCH).
\emph{arXiv}.

{[}3{]} Mahir, I. (2026). Hierarchical Cognitive Manifold Architecture
(HCMA). \emph{arXiv}.

{[}4{]} Feynman, R. (1963). The Feynman Lectures on Physics.
\emph{Addison-Wesley}.

{[}5{]} Goldstein, H. (1980). Classical Mechanics.
\emph{Addison-Wesley}.

\begin{center}\rule{0.5\linewidth}{0.5pt}\end{center}

\textbf{OMT: Think with waves. Reason with resonance. Never memorize.}

\end{document}
